% Part: normal-modal-logic
% Chapter: syntax-and-semantics
% Section: modal-validity

\documentclass[../../../include/open-logic-section]{subfiles}

\begin{document}

\olfileid[pt]{nml}{syn}{val}

\olsection{Validade}

\begin{explain}
  As !!^{formula}s que são verdadeiras em todos os modelos, isto é, em
  todos os mundos de todos os modelos, são particularmente interessantes.
  Elas representam as proposições modais verdadeiras independentemente de
  como $\Box$ e $\Diamond$ sejam interpretados, desde que a interpretação
  seja ``normal'', no sentido de ser gerada por alguma relação de
  acessibilidade sobre mundos possíveis. Chamamos tais !!{formula}s de
  \emph{válidas}. Por exemplo, $\Box(p \land q) \lif \Box p$ é válida.
  Algumas !!{formula}s que poderíamos esperar serem válidas com base na
  interpretação alética de $\Box$, como $\Box p \lif p$, contudo não são
  válidas. Parte do interesse dos modelos relacionais está no fato de que
  diferentes interpretações de $\Box$ e $\Diamond$ podem ser captadas por
  diferentes tipos de relações de acessibilidade. Isso sugere que devemos
  definir a validade não apenas em relação a \emph{todos} os modelos, mas
  em relação a todos os modelos \emph{de certo tipo}. Veremos, por exemplo,
  que $\Box p \lif p$ é verdadeira em todos os modelos nos quais cada mundo
  é acessível a partir de si mesmo, isto é, nos quais $R$ é reflexiva.
  Definir a validade em relação a classes de modelos permite formulá-lo
  sucintamente: $\Box p \lif p$ é válida na classe dos modelos reflexivos.
\end{explain}

\begin{defn}
  !!^a{formula}~$!A$ é \emph{válida} em uma classe $\mClass{C}$ de
  modelos se é verdadeira em todo modelo de~$\mClass{C}$ (isto é, em
  todo mundo de todo modelo de~$\mClass{C}$). Se $!A$ é válida
  em~$\mClass{C}$, escrevemos $\mClass{C} \Entails !A$; escrevemos
  $\Entails !A$ se $!A$ é válida na classe de \emph{todos} os modelos.
\end{defn}

\begin{prop}\ollabel{prop:subset-class}
  Se $!A$ é válida em $\mClass{C}$, também é válida em toda classe
  $\mClass{C}' \subseteq \mClass{C}$.
\end{prop}

\begin{prop}\ollabel{prop:Nec-rule}
  Se $!A$ é válida, então $\Box!A$ também é válida. 
\end{prop}

\begin{proof}
  Suponha $\Entails !A$. Para mostrar $\Entails \Box!A$, sejam
  $\mModel{M} = \tuple{W, R, V}$ um modelo e $w \in W$. Se $Rww'$,
  então $\mSat{M}{!A}[w']$, pois $!A$ é válida; portanto, também
  $\mSat{M}{\Box!A}[w]$. Como $\mModel{M}$ e $w$ eram arbitrários,
  $\Entails \Box!A$.
\end{proof}

\begin{prob}
  Mostre que as fórmulas a seguir são válidas:
  \begin{enumerate}
  \item $\Box p \lif \Box (q \lif p)$;
  \item $\Box \lnot \lfalse$;
  \item $\Box p \lif (\Box q \lif \Box p)$.
  \end{enumerate}
\end{prob}

\begin{prob}
  Mostre que $!A \lif \Box!A$ é válida na classe $\mClass{C}$ dos
  modelos $\mModel{M} = \tuple{W, R, V}$ em que $W = \{w\}$. De modo
  semelhante, mostre que $!B \lif \Box !A$ e $\Diamond !A \lif !B$
  são válidas na classe dos modelos $\mModel{M} = \tuple{W, R, V}$ em que
  $R = \emptyset$.
\end{prob}

\end{document}
